%% Question bank chapter file - generated by bank_tools.py build (do not edit while the bank is still being built:
%% it is regenerated from the per-test data in Question Bank\data). Every question is one bankq block:
%%   \begin{bankq}[id=..., source={...}, paper=.., qno=.., type=..., subtopic={...}, answer={...}, answer_src=key|solved,
%%                 working={...}, difficulty=1-10, diff_parts={K S I C F}, diff_note={...},
%%                 group=<passage id or empty>, also_in={duplicate copies in other tests}]
%% difficulty = min(10, K+S+I+C+F-1): K concepts 1-3, S steps 1-4, I insight 0-2, C computation 0-2, F format 0-1
%%   ...question LaTeX, options, figures (\bfig / \bfigi -> img/)...
%%   \end{bankq}
%% Passages shared by several questions are bankpassage blocks; their questions carry group=<passage id>.
\documentclass[10pt]{article}
\usepackage[a4paper,margin=1.8cm]{geometry}
\usepackage[T1]{fontenc}
\usepackage{amsmath,amssymb,mathtools}
\usepackage{graphicx}
\usepackage[export]{adjustbox}
\usepackage{xcolor,needspace,keyval}
\usepackage{enumitem}
\usepackage{booktabs,array,tabularx}
\usepackage[hidelinks]{hyperref}
\setlength{\parindent}{0pt}
\setlength{\parskip}{3pt}
\makeatletter
\define@key{bq}{id}{\def\bq@id{#1}}\define@key{bq}{source}{\def\bq@src{#1}}\define@key{bq}{paper}{\def\bq@paper{#1}}
\define@key{bq}{qno}{\def\bq@qno{#1}}\define@key{bq}{type}{\def\bq@type{#1}}\define@key{bq}{subtopic}{\def\bq@sub{#1}}
\define@key{bq}{answer}{\def\bq@ans{#1}}\define@key{bq}{answer_src}{\def\bq@asrc{#1}}\define@key{bq}{working}{\def\bq@work{#1}}
\define@key{bq}{group}{\def\bq@group{#1}}\define@key{bq}{also_in}{\def\bq@also{#1}}
\define@key{bq}{difficulty}{\def\bq@diff{#1}}\define@key{bq}{diff_parts}{\def\bq@dparts{#1}}\define@key{bq}{diff_note}{\def\bq@dnote{#1}}
\newenvironment{bankq}[1][]{\def\bq@work{}\def\bq@also{}\def\bq@group{}\def\bq@diff{?}\def\bq@dparts{}\def\bq@dnote{}\setkeys{bq}{#1}%
  \par\needspace{4\baselineskip}\medskip\noindent{\small\textbf{\bq@id}\hfill\color{gray}[\bq@src, Paper \bq@paper, Q\bq@qno]
  \bq@type\ $\cdot$ \bq@sub\ $\cdot$ \textbf{D\bq@diff}}\par\nopagebreak\relax}
  {\par{\small\color{teal}Answer: \textbf{\bq@ans}\ifx\bq@work\@empty\else\ (solved: \bq@work)\fi
  \ifx\bq@also\@empty\else\par\color{gray}Also in: \bq@also\fi}\par}
\define@key{bp}{id}{\def\bp@id{#1}}\define@key{bp}{source}{\def\bp@src{#1}}\define@key{bp}{questions}{\def\bp@qs{#1}}
\newenvironment{bankpassage}[1][]{\setkeys{bp}{#1}\par\needspace{4\baselineskip}\bigskip\noindent
  {\small\textbf{Passage \bp@id}\hfill\color{gray}[\bp@src; questions \bp@qs]}\par\nopagebreak\relax\itshape}{\par}
\makeatother
\newcommand{\bfig}[2]{\par\smallskip{\centering\includegraphics[width=#1pt,max width=0.62\linewidth,max height=0.32\textheight]{img/#2}\par}\smallskip}
\newcommand{\bfigi}[2]{\includegraphics[width=#1pt,max width=\linewidth,max height=0.32\textheight,valign=c]{img/#2}}
\newcommand{\bmissing}[1]{\textit{[figure missing in the source conversion: \detokenize{#1}]}}
\newenvironment{qoptions}{\begin{itemize}[leftmargin=1.8em,itemsep=1pt,topsep=2pt,label={}]}{\end{itemize}}
\newcommand{\opt}[2]{\item[(#1)] #2}
\providecommand{\fwstart}{}\providecommand{\fwend}{}\providecommand{\qfullsplit}{}
\newcolumntype{L}[1]{>{\raggedright\arraybackslash}p{#1}}
\newcolumntype{C}[1]{>{\centering\arraybackslash}p{#1}}
\newcolumntype{R}[1]{>{\raggedleft\arraybackslash}p{#1}}
\providecommand{\sectionhdr}[1]{}\providecommand{\subsectionhdr}[1]{}


\begin{document}
{\Large\bfseries KTG \& Gas Laws --- Question Bank}\par\smallskip
{\small Physics chapter P14. Questions from all converted test papers, grouped by source series. Answers: printed key where it matches the paper, otherwise solved (answer\_src=solved, with working). Exact repeats across papers appear once, with also\_in listing the other copies.}\par\medskip\hrule

\section*{TOAS papers}
\begin{bankq}[id=TOAS13-P1-PHY-S1-Q04, source={TOAS - Test 13 (24 Jul 2022)}, paper=1, qno=4, type=multi, subtopic={Kinetic theory}, answer={A,B,C,D}, answer_src=key, difficulty=3, diff_parts={K2 S1 I0 C0 F1}, diff_note={concepts: partial pressure and moles, kinetic theory | steps: mole ratio from masses and molar masses; compare pressure, number, KE per molecule and rms speed}]
A closed vessel contains a mixture of two diatomic gases A and B. Molar mass of A is 16 times that of B and mass of gas A, contained in the vessel is 2 times that of B :-
\begin{qoptions}
\opt{A}{Pressure exerted by gas B is eight times that of exerted by gas A.}
\opt{B}{Number of molecules of gas B in the cylinder is eight times that of gas A.}
\opt{C}{Average kinetic energy per molecule of gas A is equal to that of gas B.}
\opt{D}{Root mean square value of translational speed of gas B is four times that of A.}
\end{qoptions}
\end{bankq}
\begin{bankq}[id=TOAS13-P2-PHY-S1-Q10, source={TOAS - Test 13 (24 Jul 2022)}, paper=2, qno=10, type=single, subtopic={Kinetic theory}, answer={B}, answer_src=key, difficulty=1, diff_parts={K1 S1 I0 C0 F0}, diff_note={concepts: energy conservation when mixing ideal gases | steps: T = (N1 T1 + N2 T2)/(N1 + N2) = 240 K}]
Two different perfect monoatomic gases at absolute temperature 300 K and 200 K having number of molecules $2\times10^{18}$ and $3\times10^{18}$ are mixed in a container of negligible heat capacity. Temperature of mixture is :
\begin{qoptions}
\opt{A}{210 K}
\opt{B}{240 K}
\opt{C}{260 K}
\opt{D}{280 K}
\end{qoptions}
\end{bankq}
\begin{bankq}[id=TOAS14-P1-PHY-S3-Q05, source={TOAS - Test 14 (07 Aug 2022)}, paper=1, qno=5, type=numerical, subtopic={Kinetic theory}, answer={2}, answer_src=key, difficulty=7, diff_parts={K2 S2 I2 C1 F1}, diff_note={concepts: kinetic theory momentum transfer, mean molecular speed | steps: mean speed from RT/M; momentum given to the moving face from the relative velocity; drag force as density times area times v0 times mean speed}]
Consider a rectangular block of wood moving with velocity $v_0=5$m/s in a gas at temperature $27^\circ$C and mass density 2 kg/m$^3$. Velocity of block is perpendicular to it's face having area $A=1$ cm$^2$. Find drag force (in N) on the block of molar mass of gas is 10 g/mol $\left(R=\dfrac{25}{3}\dfrac{\text J}{\text{mol K}}\right)$.
\end{bankq}
\begin{bankpassage}[id=TOAS16-P2-PHY-S2-G7-8, source={TOAS - Test 16 (04 Sep 2022)}, questions={7-8}]
Consider a fixed highly evacuated cylinder in a gravity free region. Inside the cylinder piston of mass $M$ can move without friction. In each part of the cylinder, a ball of mass $m$ ($m\ll M$) is moving back and forth with very high speed $v$ and in doing so they are colliding elastically with the ends of the cylinder and with the piston. When the piston is at the center, it is in equilibrium as shown in figure. Take $M=0.2$ kg, $m=10^{-3}$ kg $L=2$ m, $v=\pi$ m/s
\bfig{134.7}{TOAS16_p2_phy_s2_para7_8_1.png}
\end{bankpassage}
\begin{bankq}[id=TOAS16-P2-PHY-S2-Q07, source={TOAS - Test 16 (04 Sep 2022)}, paper=2, qno=7, type=single, subtopic={Kinetic theory}, answer={B}, answer_src=key, difficulty=6, diff_parts={K2 S2 I2 C1 F0}, diff_note={concepts: elastic collisions of a ball with a piston, oscillation | steps: force from momentum transfer; effective restoring force for small displacement; period}, group=TOAS16-P2-PHY-S2-G7-8]
Now the piston is shifted slightly towards one of the ends of cylinder with a negligible speed and released. Time period of oscillations of piston is :-
\begin{qoptions}
\opt{A}{1.5}
\opt{B}{20}
\opt{C}{3.2}
\opt{D}{45}
\end{qoptions}
\end{bankq}
\begin{bankq}[id=TOAS16-P2-PHY-S2-Q08, source={TOAS - Test 16 (04 Sep 2022)}, paper=2, qno=8, type=single, subtopic={Kinetic theory}, answer={B}, answer_src=key, difficulty=2, diff_parts={K1 S1 I1 C0 F0}, diff_note={concepts: pressure balance by molecular collisions | steps: pressure proportional to m v\^{}2/length on each side; lengths from the equal pressure}, group=TOAS16-P2-PHY-S2-G7-8]
If the ball on right side has double mass as well as speed of left side ball then length of right section of cylinder in steady state is :-
\begin{qoptions}
\opt{A}{$\dfrac{4L}5$}
\opt{B}{$\dfrac{8L}9$}
\opt{C}{$\dfrac{7L}8$}
\opt{D}{$\dfrac{3L}4$}
\end{qoptions}
\end{bankq}
\end{document}
